\documentclass[book.tex]{subfiles}

\begin{document}

\chapter{Gaussian Processes}
\label{chap:gaussian_processes}
\noindent\rule{11cm}{0.4pt} \\
\textbf{Prerequisites:} \autoref{chap:calculus}, \autoref{chap:multidimensional}, \autoref{chap:ml_basics}  \\
\textbf{Difficulty Level:} ** \\
\noindent\rule{11cm}{0.4pt}
\vspace{5mm}

Gaussian processes provide a powerful class of methods for modeling time series data that competes broadly with RNNs but offers much tighter theoretical control 

\section{Intuition}

The core idea of a Gaussian process is that a dataset
\begin{align}
\mathcal{D} &= \{(x_1, y_1),\dotsc, (x_n, y_n)\}
\end{align}
corresponds to a family of functions which pass through these points. By using as assumption of Gaussianity, we can place a reasonable prior distribution on the family of all functions which pass through the points in $\mathcal{D}$. In order for this prior to make sense at points not in $\{x_i\}$, we need to make a choice of a kernel function $K$ that tells us how to compute the probability of values $y$ at some point $x$ not in the training distribution.

In order to do prediction at $x$, we have to compute the conditional Gaussian distribution given $\mathcal{D}$ and $K$. This computation can become expensive, so there are number of methods to compute approximate Gaussian process regressors.
\end{document}